\scanpage{13}
\mspage{11}
Il convient de déterminer l'hom.
\[
I'_t\simeq\prod_{\ell\ne p}\mathbb Z_\ell(1)\longrightarrow\inserted{I_t\simeq}\prod_{\ell\ne p}\mathbb Z_\ell(1)\tag*{(3.8)}
\]
où on identifie, grâce à l'inclusion \(\overline K\subset\overline K'\), les \deleted{\unsure{deux|les}} groupes \(\mathbb Z_\ell(1)\) définis à l'aide de \(\overline K\) ou de \(\overline K'\) (savoir \(T_\ell(\overline K^*)\) et \(T_\ell(\overline K'{}^*)\)).

Pour ceci, on regarde les sous-corps (pour \((n,p)=1\))
\[
\left\{\begin{aligned}
K_n&=\widetilde K(\unsure{\Theta|\theta}_n)\subset\overline K_t&\text{où}\quad&\Theta_n^n=t,\\
K'_n&=\widetilde K'(\Theta'_n)\subset\overline K'_t&\text{où}\quad&\Theta'_n{}^n=t'
\end{aligned}\right.\tag*{(3.9)}
\]
et on regarde comment
\[
\operatorname{Gal}(K'_n/\widetilde K')\simeq\mu_n(\widetilde K')\xleftarrow{\sim}\mu_n(\widetilde K)
\]
(gpe des racines \(n\)ièmes \(\varepsilon\) de 1 ds \(\widetilde K\)) opère sur \(K_n\). Si \(\varepsilon_{K'}\) désigne l'opération sur \(K'_n\) définie par \(\varepsilon\), donc
\[
\varepsilon_{K'_n}(\Theta'_n)=\varepsilon\Theta'_n
\]
on trouve, si
\[
t=\deleted{\unsure{t'}}u't'{}^e\quad\inserted{(e\in\mathbb N,\ u'\in V^*)}\quad
\text{d'où }\deleted{\unsure{=}}\ \boxed{\begin{aligned}&\deleted{\widetilde{\mathfrak m}\widetilde V=\widetilde{\mathfrak m}'\unsure{V'}}\\&\mathfrak m V=\mathfrak m'{}^e\end{aligned}}\tag*{(3.10)}
\]
que
\[
\Theta_n^n=u'(\Theta'_n{}^e)^n
\]
donc, si on pose
\[
\Theta_n=\lambda\Theta'_n{}^e,\quad\inserted{(\lambda\in K'_n)}\tag*{(3.11)}
\]
on aura
\[
\lambda^n=u',\tag*{\deleted{\unsure{(3.12)}}}
\]
donc \deleted{\unsure{reste|d'où}} \(\lambda\in\widetilde V'\) \deleted{\unsure{et}} donc \(\varepsilon_{K'}(\lambda)=\lambda\) et on trouve
\[
\varepsilon_{K'_n}(\Theta_n)=\underset{\substack{\shortparallel\\\lambda}}{\varepsilon_{K'_n}(\lambda)}\varepsilon_{K'_n}(\Theta'_n{}^{\deleted{\unsure{e}}})^e=\lambda(\varepsilon\Theta'_n)^e\overset{\inserted{\text{par (3.10)}}}{=}\varepsilon^e\Theta_n
\]
% archive sheet number: 12
